%%%%%
%
%% Dans ce fichier, les �ventuels accents sont cod�s
%% en ISO-8859-15 (aka latin9)
%% La ligne suivante donne trois exemples 
%% e accent aigu : �, e accent grave : �, e-dans-l'o : �
\section{Rapport de m�lange saturant}

\itemstyle{Op\'erateur~:} \operfmt{RapM\_satur}
\index{RapM\_satur}

%%%%%OBJET%%%%
\itemstyle{Objet~:}
Calcule le rapport de m�lange saturant, c'est � dire la valeur
maximale que le rapport de m�lange en vapeur d'eau d'une parcelle
d'air peut atteindre pour une pression et une temp�rature donn�e, 
au dessus d'une surface plane d'eau liquide pour les 
temp�rature positives (en �C) et au-dessus de la glace lorsque
la temp�rature est n�gative (toujours en �C).

%%%%%ARGUMENTS%%%%
\bigskip
%% Dans ce fichier, les �ventuels accents sont cod�s
%% en ISO-8859-15 (aka latin9)
%% La ligne suivante donne trois exemples 
%% e accent aigu : �, e accent grave : �, e-dans-l'o : �
%%%%%ENTREES%%%%
\begin{tabarg}
\multicolumn{3}{l}{\titretab {Entr\'ees}}\\
\hline
$T_{s}$ & Vecteur, \texttt{float} & Temp�rature d'�tat ($K$ ou $^\circ C$)\\
$P_{s}$ & Vecteur, \texttt{float} & Pression ($hPa$)\\
\end{tabarg}

\medskip
%%%%%SORTIES%%%%
\begin{tabarg}
\multicolumn{3}{l}{\titretab {Sorties}}\\
\hline
$RapM_{satur}$ & Vecteur, \texttt{float} & Rapport de m�lange saturant ($g/kg$)\\
\end{tabarg}


%%%%%FORMULE%%%%
\itemstyle{Formule}

%% Formule de l'operateur : RapM_satur
%% Dans ce fichier, les �ventuels accents sont cod�s
%% en ISO-8859-15 (aka latin9)
%% La ligne suivante donne trois exemples 
%% e accent aigu : �, e accent grave : �, e-dans-l'o : �
\[
e = 
\begin{cases}
e_{w}(T_{s}) \text{ si } T_{s} > 0�C \\
e_{wi}(T_{s}) \text{ sinon }
\end{cases}
\]
\[
RapM_{satur} = 1000 . \frac{R_a}{R_v} . \frac{e}{(P_{s} - e)}
\]



%%%%%AUTEUR%%%%
\itemstyle{Auteur~:}
CNRM/GMEI/TRAMM

%%%%%REFERENCE%%%%
\itemstyle{R�f�rences~:}
